\section{Introduction}
Building energy controllers must balance cost, carbon, peak demand, discomfort, degradation, ramping, and resilience. For each of three buildings, the controller receives seven non-negative weights; a candidate therefore contains 21 values, with each building's seven values normalized to sum to one in our experiments.

The project objective is to minimize the district ramping score returned by a deterministic black-box evaluator:
\begin{equation}
R(W)=\frac{1}{8}\sum_{w=1}^{8}\sum_{t=2}^{168}|E_{w,t}(W)-E_{w,t-1}(W)|,
\end{equation}
where $E_{w,t}$ denotes district net load at hour $t$ of week $w$. Lower $R$ is better. The online stage permits at most 20 sequential evaluator calls, each containing at most 50 candidate weight sets. Before the first call, the policy may use only the 2,048-row training table; subsequent calls may append measured results and retrain the surrogate. This makes duplicate prevention, restart safety, and provenance as important as the choice of model.

Our work separates three questions: whether a surrogate ranks low-score regions accurately, how an optimizer turns its predictions into a diverse batch, and whether the complete policy remains reproducible under the evaluator budget. The report compares these components with the frozen-v3.2 protocol and then records the online result separately from offline predictions.
